% META: TODO.1

\chapternoindex{Аннотация}

\textbf{Объём работы:} (TODO) страниц, включая (TODO) рисунков, (TODO) таблиц, (TODO) приложений. Библиографический список содержит (TODO) наименований.

\textbf{Ключевые слова:} машинное обучение, нейронные сети, CAD, КОМПАС-3D, обратный инжиниринг, история построения, сегментация 3D-моделей, MeshCNN, STEP, STL.

\textbf{Объект исследования:} процесс восстановления истории построения трёхмерных параметрических моделей в системе автоматизированного проектирования.

\textbf{Предмет исследования:} методы машинного обучения (глубокие нейронные сети для анализа трёхмерных сеток), применяемые для семантической сегментации геометрии и распознавания формообразующих операций, а также способы их интеграции с детерминированными алгоритмами CAD-ядра.

\textbf{Цель работы:} разработка программного модуля «НейроКОМПАС-3D», интегрируемого со средой КОМПАС-3D, который на основе анализа геометрии модели позволяет реконструировать последовательность породивших её операций.

\textbf{Методы исследования:} системный анализ, методы глубокого обучения (архитектура MeshCNN для классификации и сегментации), методы вычислительной геометрии (триангуляция, нормализация сеток, фильтрация), экспериментальное тестирование с количественным анализом метрик.

\textbf{Краткое содержание.} В работе проведён анализ предметной области, выявлены ограничения подхода прямого предсказания CAD-команд (высокая чувствительность регрессии параметров к шумам и системе координат). Обоснован и реализован гибридный подход, сочетающий нейросетевую сегментацию геометрии для распознавания следов операций (выдавливание, вырезание, скругление, фаска) с последующим детерминированным вычислением параметров средствами CAD-ядра.

Разработан конвейер подготовки датасета, позволяющий из сырых данных (истории построения в виде STEP файлов каждого шага + JSON-описание названий операций) автоматически генерировать размеченные по рёбрам OBJ-файлы. Исходный массив из (TODO)82 418 моделей после очистки и фильтрации сокращён до рабочего набора в (TODO)21 464 модели. На основе архитектуры MeshCNN обучены классификаторы и сегментаторы для четырёх основных операций: bossExtrusion, cutExtrusion, fillet, chamfer.

Экспериментально подтверждена состоятельность подхода согласно метрикам, измеренным на тестовой выборке после обучения. Точность классификации (accuracy, равное отношению числа верно классифицированных рёбер к общему числу рёбер) достигла: bossExtrusion - (TODO)\%, cutExtrusion - (TODO)\%, fillet - (TODO)90,6\%, chamfer - (TODO)\%. Метрика IoU: bossExtrusion - (TODO)\%, cutExtrusion - (TODO)\%, fillet - (TODO)\%, chamfer - (TODO)\%. Точность сегментации (accuracy): bossExtrusion - (TODO)\%, cutExtrusion - (TODO)\%, fillet - (TODO)\%, chamfer - (TODO)\%. Доля моделей, в которых минимум 90\% рёбер сегментированы верно, составила: bossExtrusion - (TODO)\%, cutExtrusion - (TODO)\%, fillet - (TODO)\%, chamfer - (TODO)\%.

Создан исполняемый модуль main.exe, реализующий полный цикл инференса: загрузку STEP / STL, адаптивную триангуляцию с бинарным поиском для ограничения числа треугольников (не более 1500), нормализацию сетки до фиксированного числа рёбер (2250), вызов нейросетевых моделей в автоматическом (перебор операций с исключением отвергнутых гипотез) или ручном режиме, постобработку (фильтрация треугольников по площади, выделение связного компонента) и подготовку выходных данных (STL-файл выделенной геометрии, JSON-файл с типом операции и уверенностью корректного её определения).

Модуль интегрируется с КОМПАС-3D через API, позволяя выполнять упрощение модели путём пошагового удаления распознанных операций. Разработанная архитектура допускает независимое дообучение классификаторов и сегментаторов, а также добавление новых типов операций без переучивания ранее добавленных, что даёт модульность и возможность улучшения результатов только для необходимых типов операций.

\textbf{Структура работы.} В первой главе проведён анализ предметной области, формализована задача, обоснован выбор гибридного подхода и нейросетевой архитектуры MeshCNN. Во второй главе описана архитектура решения: конвейер подготовки данных, схемы обучения и инференса, требования к форматам представления, а также приведены результаты экспериментального обучения и тестирования, а также сформулированы предложения по дальнейшему развитию (увеличение датасета, поддержка дополнительных операций, улучшение постобработки). (TODO)В приложениях вынесена техническая документация по установке, настройке и эксплуатации, а также примеры работы модуля на типовых деталях.

\textbf{Практическая значимость.} Разработанный модуль может быть использован для упрощения сложных 3D-моделей, восстановления параметрической истории импортированных моделей и обучения инженеров-конструкторов. Результаты работы внедрены в опытную эксплуатацию в рамках договора с ООО «АСКОН-Системы проектирования».

